\documentclass[10pt,twocolumn]{article}
\usepackage[margin=1.8cm]{geometry}
\usepackage{amsmath}
\usepackage{booktabs}
\usepackage{hyperref}
\usepackage{parskip}
\setlength{\parskip}{4pt}

\title{\textbf{Knowledge Tracing with Vision-Language Models}\\
\large Two Approaches to Student Performance Prediction}
\author{Guglielmo Montone}
\date{}

\begin{document}
\maketitle
\thispagestyle{empty}

\section*{Problem Statement}

Knowledge Tracing (KT) is the task of predicting whether a student will answer the next exercise
correctly, given their history of past attempts. This work explores two experimental pipelines
trained and evaluated on the \textbf{MIAAM dataset}, which contains 38,519 students with an
average of 168 exercises each. Both pipelines leverage \textbf{SmolVLM-256M-Instruct} as the
multimodal backbone and share a common data representation: exercises are described by text
and a screenshot, while attempts are described by text (answer given, duration, timestamp,
and correctness label).

The first experiment is intended as a \textbf{baseline}: feeding raw multimodal content directly
to a VLM is the most straightforward approach to the problem, and it helps characterise the
limitations of this strategy on the MIAAM dataset. The second experiment proposes an alternative
design in which the VLM is used to encode each exercise and each attempt \emph{independently}
into a fixed-size vector; these representations are then consumed by a Transformer architecture,
which can reason over much longer student histories.

\section*{Experiment 1 --- Raw Multimodal Input}

The first approach feeds \emph{raw} multimodal content directly to the VLM at training time. 
For a student with history $t$, the model receives a concatenated sequence of exercise 
screenshots and text, together with text descriptions of past attempts.

\textbf{Objective.} Establish a baseline by fine-tuning a VLM either end-to-end or 
partially, directly on the raw multimodal input.

Two training regimes are supported:
\begin{itemize}
  \item \textbf{Probe} — the VLM backbone is fully frozen; only a linear classification head appended
   to the pooled hidden state is trained. This is fast and parameter-efficient.
  \item \textbf{LoRA fine-tuning} — lightweight LoRA adapters are injected into the attention 
  and MLP projections of the backbone and trained end-to-end alongside the head, at a lower 
  learning rate than the probe.
\end{itemize}

\textbf{Key limitations.} The 256M-parameter VLM has a small context window, which restricts each
input to only 4--5 past attempts regardless of the student's actual history length. Since students
average 168 exercises, the vast majority of each student's learning trajectory is discarded at
every training step, inherently limiting what the model can infer about long-term knowledge state.

A second limitation is \textbf{training speed}. Even in probe mode, where the backbone is frozen and
no gradients flow through it, a full VLM forward pass must run on every batch to produce the hidden
states consumed by the classification head. Processing raw screenshots through the vision encoder at
every step is expensive, and this cost is paid repeatedly across all epochs. By contrast, Experiment~2
pays the VLM inference cost only once during preprocessing: embeddings are cached to disk and the
training loop never touches the VLM again, making each training step orders of magnitude cheaper.

\section*{Experiment 2 --- Pre-computed Embeddings}

The second approach decouples the VLM from the training loop entirely. Each exercise and each
attempt is encoded \emph{once} by a frozen VLM into a 576-dimensional embedding vector. A
lightweight Transformer then learns to reason over sequences of these cached embeddings.

\textbf{Objectives.} (1) Validate the advantage of pre-computing exercise and attempt embeddings 
from a computational perspective, expecting faster training that can run on smaller GPUs. (2) Validate 
the advantage in terms of prediction performances of feeding the model a longer sequence of 
student history.

\textbf{Input sequence.} For a student who has completed $t$ exercises, the model 
receives an interleaved sequence of $2t+1$ vectors:
\[
[\,e_0,\,a_0,\,e_1,\,a_1,\,\ldots,\,e_{t-1},\,a_{t-1},\,e_t\,]
\]
where $e_i = \mathrm{VLM}(\text{exercise}_i)$ and $a_i = \mathrm{VLM}(\text{attempt}_i)$. The final 
token $e_t$ has no corresponding attempt yet; the prediction head reads the Transformer's hidden 
state at that position and produces a scalar logit for binary classification.

\textbf{Architecture.} The \texttt{KTTransformer} consists of (i) a linear projection from the 576-d 
input space to the model's internal dimension $d_\text{model}$, (ii) learnable positional 
embeddings, (iii) a stack of standard Transformer encoder layers (bidirectional attention, no causal mask), 
and (iv) a two-layer MLP head. Sequences are zero-padded to the batch maximum and a boolean attention
mask prevents the model from attending to padding tokens.

\textbf{Key advantage.} Because embeddings are pre-computed and compact, the Transformer can process 
histories of up to \textbf{512 past exercises}, covering most students in full rather than truncating 
their history.

\section*{Results}

No complete results was found for Experiment~1. A single training epoch on the available GPU takes 3--4 
hours due to the cost of running the full VLM forward pass at every step, making a full training run 
impractical within the available time budget.

The embedded approach reaches \textbf{83\% validation accuracy} after a short training session, 
with no hyperparameter tuning. This is a strong baseline given the simplicity of the 
encoder (256M parameters) and the absence of any hand-crafted features.

\section*{Conclusion}

Experiment 1 lets the VLM jointly encode each exercise and reason over the full history 
end-to-end, so it can interpret an exercise in light of what just happened (e.g. right after 
a failure vs. a success), at the cost of higher capacity demands and the risk of entangling 
exercise content with spurious features of any one student's trajectory. Experiment 2 factorizes 
the problem — a frozen encoder produces consistent, context-free exercise embeddings that a 
dedicated downstream transformer contextualizes over the sequence — which is more parameter-efficient
and avoids baking history-specific noise into the content representation, but loses the ability 
to fuse context at encoding time if the embedding discards the relevant detail.

Regarding Experiment~1, it would still be interesting to run a complete training on this
setup, potentially using a more powerful model with a larger context window to better
capture student history.

Experiment~2 appears very promising, especially when one of the main desiderata of the
project is to operate with small VLMs. It offers clear advantages in training speed and
in the ability to feed the model a much longer student history. The natural next step
here is to try a slightly larger backbone to produce richer embeddings, while retaining
the decoupled, computationally efficient design.

\end{document}
